% TOP_PROBLEM: {"id": 30006867, "title": "Discontinuity of the q-State Potts Phase Transition on Z^d for d >= 3 and q >= 3", "category_id": 16, "category": {"id": 16, "name": "physics", "display_name": "Mathematical Physics", "description": "Problems at the intersection of mathematics and physics.", "slug": "physics", "order_index": 16, "created_at": "2026-07-31T15:26:25.671Z"}, "status": "open"}
% FIRST_POSTED: 2026-09-27
% !TeX program = lualatex
% ATTEMPT_STATUS: unresolved
% Model: GPT-6 Astra Ultra; continuation dated 27 September 2026.
\documentclass[11pt]{article}
\usepackage[margin=25mm]{geometry}
\usepackage{fontspec}
\setmainfont{Latin Modern Roman}
\usepackage{amsmath,amssymb,mathtools}
\usepackage{unicode-math}
\setmathfont{Latin Modern Math}
\usepackage{xurl,hyperref}
\hypersetup{colorlinks=true,urlcolor=blue}
\setcounter{secnumdepth}{0}
\setlength{\parindent}{0pt}
\setlength{\parskip}{5pt}
\setlength{\emergencystretch}{3em}
\begin{document}
\section*{\texorpdfstring{Discontinuity of the q-State Potts Phase Transition on Z\textasciicircum{}d for d >= 3 and q >= 3}{Discontinuity of the q-State Potts Phase Transition on Z\textasciicircum{}d for d >= 3 and q >= 3}}\label{30006867}
\subsection{Definitions and mathematical statement}
In the nearest-neighbor ferromagnetic Potts model, spins $\sigma_x\in\{1,\ldots,q\}$ on ${\mathbb{Z}}^d$ have finite-volume weights proportional to
\[
 \exp\!\left(\beta\sum_{\langle x,y\rangle}
 \mathbf1_{\{\sigma_x=\sigma_y\}}\right),
\]
with fixed boundary spins used to select an ordered state. Write $m_q(\beta)$ for its excess probability of spin $1$ at the origin above $1/q$, in the infinite-volume limit, and let $\beta_c$ be the threshold for positive order. The order-parameter discontinuity formulation asks $m_q(\beta_c)>0$ for every $d\ge3$, $q\ge3$. Equivalently in the random-cluster convention, the wired critical measure has positive probability of an infinite cluster, with $p=1-e^{-\beta}$. The catalogue calls the transition first-order but does not specify a latent-heat normalization; the displayed statement specifies the order/disorder discontinuity convention rather than adding a separate quantitative energy-jump demand.
\par\smallskip\textit{Formulation and concept references:} \href{https://ar5iv.labs.arxiv.org/html/1707.00520}{[S1]}.

\subsection{Short English statement}
Does every nearest-neighbor Potts model with at least three spin states in dimension at least three acquire long-range order discontinuously at its transition?
\subsection{Research attempt}
\textbf{The exact discontinuity under discussion.}
The target is positive order at the critical point in the wired state.  It is
not a numerical lower bound for latent heat.  The higher-dimensional discussion
and Conjecture 2 in [S1] distinguish the full nearest-neighbor assertion from
the established large-$q$ and sufficiently-high-dimensional regimes.  Those
regimes do not cover every pair $(d,q)$ in the statement.  This attempt derives
the cluster formulation, proves elementary noncritical estimates, and tests a
tree recursion which displays the expected discontinuous mechanism but does
not transfer to the lattice.

\textbf{Exact finite-volume spin--cluster correspondence.}
Let $G=(V,E)$ be finite, with its boundary vertices wired together and assigned
spin $1$.  Put $p=1-e^{-\beta}$.  For each edge,
\[
 e^{\beta\mathbf1_{\{\sigma_x=\sigma_y\}}}
 =e^\beta\bigl((1-p)+p\mathbf1_{\{\sigma_x=\sigma_y\}}\bigr).
\]
Expanding the product over edges and then summing spins gives the
random-cluster weight
\[
 p^{|\omega|}(1-p)^{|E|-|\omega|}q^{k_0(\omega)},
\]
apart from the constant $e^{\beta|E|}$.  Here $k_0$ counts open components not
connected to the wired boundary.  Given $\omega$, assign spin $1$ to the
boundary cluster and independent uniform spins to the other clusters.
If the origin is not boundary-connected its spin is uniform; otherwise it is
$1$.  Consequently
\[
 \mathbb P(\sigma_0=1)-\frac1q
 =\left(1-\frac1q\right)\phi^w_{G,p,q}(0\leftrightarrow\partial G).
\]
The same identity holds in the wired infinite-volume limit.  In particular,
the normalization in this entry is
$m_q(\beta)=(1-1/q)\theta^w(p)$, not $\theta^w(p)$ itself.

\textbf{Conditional edge probabilities and elementary bounds.}
Fix all edges except $e$.  If its endpoints are already connected, opening
$e$ does not change $k_0$ and its conditional open probability is $p$.
If they are not connected, opening merges two components, at least one of
which is not the wired component, and reduces $k_0$ by one.  The probability
is then
\[
 r=\frac{p}{p+q(1-p)}.
\]
Thus every single-edge conditional probability lies in $[r,p]$.
Revealing edges in any fixed order and using independent uniform random
variables couples the configuration between Bernoulli configurations of
parameters $r$ and $p$: each conditional probability given the revealed
prefix is an average of the full conditional probabilities and remains in
that interval.  This proves the finite-volume stochastic bounds directly.
Limits preserve inequalities for increasing cylinder events, hence give the
corresponding infinite-volume comparisons.

For a self-avoiding path with $n$ edges, its all-open probability is at most
$p^n$.  There are at most $2d(2d-1)^{n-1}$ such paths starting at the origin.
An infinite open cluster contains paths of every length, so
\[
 \theta^w(p)\leq 2d(2d-1)^{n-1}p^n,
 \qquad p<\frac1{2d-1}\ \Longrightarrow\ \theta^w(p)=0.
\]
The argument works for every $q\geq1$ and makes no critical-point claim.

For completeness, a low-temperature positivity bound can be proved by
restricting the dominated Bernoulli configuration to a coordinate copy of
$\mathbb Z^2$.  If the origin's Bernoulli cluster is finite, a closed-edge
dual circuit surrounds it.  The number of length-$n$ candidates is bounded
by $4n3^n$: choose an intersection with a fixed positive coordinate ray and
then a nonbacktracking walk.  The crude bound suffices.  Therefore
\[
 \mathbb P_r(0\not\leftrightarrow\infty)
 \leq\sum_{n\geq4}4n\,[3(1-r)]^n.
\]
For $r$ sufficiently close to one this sum is below one, proving
$\theta^w(p)>0$.  Since $r\uparrow1$ as $p\uparrow1$ for fixed finite $q$,
these estimates bracket a nontrivial threshold.  They leave a large interval
between the explicit high- and low-temperature conditions, and do not give
a positive lower bound at its actual critical endpoint.

\textbf{A useful limit-order check.}
Let $\Lambda_N$ be increasing boxes and
$\theta_N(p)=\phi^w_{\Lambda_N,p,q}(0\leftrightarrow\partial\Lambda_N)$.
For fixed $N$ this is a ratio of finite polynomials with positive denominator
on $0<p<1$, hence continuous there.  The spatial Markov property and the
maximality of wired boundary conditions imply that $\theta_N$ decreases with
$N$.  Its limit is $\theta^w(p)$.  To check the latter point rather than assume
an exchange of limits, fix a smaller box $\Lambda_R$.  For $N\geq R$,
\[
 \theta_N(p)\leq\phi^w_{\Lambda_N,p,q}(0\leftrightarrow\partial\Lambda_R).
\]
Taking $N\to\infty$ first gives the infinite-volume probability on the right;
then $R\to\infty$ gives $\theta^w(p)$.  The reverse inequality follows from
the wired-domain comparison.  The standard monotone construction of the
wired measure underlies these comparisons.

Thus $\theta^w(p)=\inf_N\theta_N(p)$ is upper semicontinuous.  It is also
nondecreasing in $p$, by the random-cluster monotonicity for $q\geq1$.
These two facts imply right continuity at interior parameters:
\[
 \theta^w(p)=\lim_{s\downarrow p}\theta^w(s).
\]
They do not imply left continuity.  In fact the conjecture asks for a jump
from the zero left-hand regime to a positive value at $p_c$.  Finite-volume
analyticity does not prohibit such a jump after the volume limit.
Equally, right continuity does not prove a jump: the function
$\max\{p-p_c,0\}$ has all these properties and equals zero at $p_c$.
This elementary model is a countercheck on the attempted inference, not a
candidate Potts order parameter.

\textbf{An exactly solvable recursion showing a discontinuous branch.}
On a rooted tree with $b$ children per vertex, let the incoming branch
partition functions favor color $1$ by a ratio $z$ over every other color.
Writing $E=e^\beta$, summation over each child's spin yields the recursion
\[
 z\longmapsto\left(\frac{Ez+q-1}{z+E+q-2}\right)^b.
\]
The symmetric solution $z=1$ is always fixed.  The derivative of this map
there equals $b(E-1)/(E+q-1)$, so its linear stability changes at
$E=(b+q-1)/(b-1)$.

For the concrete choice $b=2$, $q=3$, put $y=\sqrt z>0$.  The fixed-point
equation becomes
\[
 y^3-Ey^2+(E+1)y-2
 =(y-1)\bigl(y^2+(1-E)y+2\bigr)=0.
\]
The nonsymmetric positive roots first appear at
$E_*=1+2\sqrt2$, when $y=\sqrt2$ and $z=2$.
They therefore appear a definite distance from $z=1$, while the latter's
linear derivative is still less than one because $E_*<4$.
This gives an explicit coexistence-of-fixed-points mechanism rather than a
continuous bifurcation from the symmetric branch.  At $E=4$ the two extra
roots are $y=1,2$, agreeing with the computed loss of linear stability.

These equations describe a tree message recursion, not the nearest-neighbor
lattice model.  Even on the tree, selecting a thermodynamic phase requires
specifying boundary conditions and the relevant free energy; the mere birth
of two fixed points is not a proof that this value is the transition point.
On $\mathbb Z^d$, loops correlate incoming branches, so multiplying $b$
identical independent contributions is not an exact calculation.  Replacing
$b$ by $2d-1$ cannot repair this loss of independence.

\textbf{Where a contour attack must gain a uniform estimate.}
The cluster bounds suggest trying to separate an ordered ensemble from a
disordered one and show that interfaces between them are costly.  An
estimate that works only when $p$ is very close to one proves low-temperature
order, which is already supplied above.  To prove the target it must remain
quantitatively valid at the actual coexistence parameter, without assuming
that $q$ or $d$ is large.  This requires enough interface suppression to keep
the wired infinite cluster density bounded away from zero as
$p\downarrow p_c$.  Neither the single-edge comparison nor the tree recursion
establishes such suppression for $(d,q)=(3,3)$, the smallest requested case.
Numerical free-energy crossings also do not bound a thermodynamic-limit
order parameter without uniform errors and a controlled volume limit.

\textbf{Verification and outcome.}
The diagnostic enumerates all spin and bond configurations of a small wired
graph at rational $e^\beta$, checking the expansion and the exact factor
$1-1/q$.  It checks the tree polynomial factorization with exact integer
coefficients and evaluates the crude Peierls sum at one explicit high-density
parameter.  These are finite identities and a numerical bound, not evidence
for a uniform critical jump.  The first missing step remains a strictly
positive critical-density estimate valid for every requested $d$ and $q$.
\subsection{Sources}
\begin{itemize}
\item[S1]H. Duminil-Copin, 'Lectures on the Ising and Potts models on the hypercubic lattice', arXiv:1707.00520 (2017), in Random Graphs, Phase Transitions and the Gaussian Free Field, Springer PIMS 123 (2020).
\url{https://ar5iv.labs.arxiv.org/html/1707.00520}.
\textit{Formulation source.}
\end{itemize}
\end{document}
